\documentclass[11pt]{article}

% Page layout
\usepackage[a4paper, margin=1in]{geometry}

% Math packages
\usepackage{amsmath, amssymb, amsfonts}

% Typography
\usepackage{lmodern}
\usepackage{microtype}
\usepackage{setspace}
\setstretch{1.1}

% Lists
\usepackage{enumitem}
\setlist{noitemsep}

% Tables
\usepackage{array}
\usepackage{booktabs}

% Colors and boxes
\usepackage{xcolor}
\usepackage[most]{tcolorbox}

% TikZ for diagrams
\usepackage{tikz}

% Hyperlinks
\usepackage[hidelinks]{hyperref}

% Section styling
\usepackage{titlesec}

\titleformat{\section}
{\Large\bfseries}{}
{0em}{}

\titleformat{\subsection}
{\large\bfseries\color{black!80}}{}{0em}{}

% Paragraph formatting
\setlength{\parindent}{0pt}
\setlength{\parskip}{6pt}

% Separator
\newcommand{\sepline}{
    \vspace{0.5em}
    \hrule
    \vspace{1em}
}

% Colors
\definecolor{myblue}{HTML}{D6EAF8}
\definecolor{myred}{HTML}{FADBD8}
\definecolor{mypurple}{HTML}{E8DAEF}

% Boxes
\newtcolorbox{theorybox}{
    colback=mypurple,
    colframe=purple!70!black,
    breakable
}

\newtcolorbox{remarkbox}{
    colback=myblue,
    colframe=blue!60!black,
    breakable
}

\newtcolorbox{problembox}{
    colback=myred,
    colframe=red!70!black,
    breakable
}

% Document
\begin{document}

\begin{titlepage}
\centering
\vspace*{3cm}

{\Huge \bfseries Cartesian Products and Ordered Pairs \par}

\vspace{1cm}

{\large Lostless1907 \par}

\vfill

\sepline

{\large \textit{A supposedly delicious introduction to pairing and dimension in sets}}

\end{titlepage}

\section*{\textsc{Preface}}
\sepline

Sets allow us to collect objects.

However, many mathematical situations require us to relate objects while preserving order.

This leads to the idea of ordered pairs and Cartesian products.

\section*{\textsc{Cartesian Products and Ordered Pairs}}
\sepline

\begin{theorybox}

\textbf{Why Order Matters}

Consider:
\[
(a,b)
\]

We require:
\[
(a,b) \neq (b,a)
\]

Thus, unlike sets, order is essential.

\textbf{Defining Ordered Pairs}

We define $(a,b)$ such that:
\[
(a,b) = (c,d) \iff a = c \text{ and } b = d
\]

This property uniquely determines ordered pairs.

\textbf{Construction}

We define:
\[
(a,b) := \{\{a\}, \{a,b\}\}
\]

This ensures the defining property holds.

\textbf{Cartesian Product}

Given sets $A$ and $B$:
\[
A \times B = \{(a,b) \mid a \in A,\; b \in B\}
\]

\textbf{Size}

If $|A| = m$ and $|B| = n$:
\[
|A \times B| = mn
\]

\textbf{Order Matters Again}

\[
A \times B \neq B \times A
\]

\end{theorybox}

\subsection*{Examples}

\begin{theorybox}

Let:
\[
A = \{1,2\}, \quad B = \{x,y\}
\]

Then:
\[
A \times B =
\{(1,x),(1,y),(2,x),(2,y)\}
\]

\[
B \times A =
\{(x,1),(x,2),(y,1),(y,2)\}
\]

\end{theorybox}

\subsection*{Geometric Interpretation}

If $A, B \subseteq \R$, then each $(a,b)$ is a point.

\begin{center}
\begin{tikzpicture}[scale=1]

\draw[->] (-1,0)--(4,0) node[right]{$x$};
\draw[->] (0,-1)--(0,4) node[above]{$y$};

\fill (1,2) circle (2pt) node[right]{$(1,2)$};
\fill (2,1) circle (2pt) node[right]{$(2,1)$};

\end{tikzpicture}
\end{center}

Thus:
\[
\R \times \R = \text{the plane}
\]

\begin{remarkbox}

A Cartesian product does not introduce new elements.

It introduces structure by pairing existing elements with order.

\end{remarkbox}

\begin{remarkbox}

Cartesian products increase dimension:

\begin{itemize}
\item $\R$ is a line
\item $\R \times \R$ is a plane
\item $\R^3$ is space
\end{itemize}

\end{remarkbox}

\begin{problembox}

\textbf{Problems}

\begin{enumerate}

\item Show that $(1,2) \neq (2,1)$.

\item If $(a,b) = (c,d)$, prove $a=c$ and $b=d$.

\item Construct $\{1,2\} \times \{a,b,c\}$.

\item Find $|\{1,2,3\} \times \{x,y\}|$.

\item Show:
\[
A \times B = \emptyset \iff A = \emptyset \text{ or } B = \emptyset
\]

\item When is $A \times B = B \times A$?

\item Show:
\[
(A \times B) \cap (C \times D) = (A \cap C) \times (B \cap D)
\]

\item Describe $A \times A$.

\item If $A = \R$, describe $A \times A$.

\item Challenge:

How many subsets does $\{1,2\} \times \{x,y\}$ have?

\end{enumerate}

\end{problembox}

\end{document}
